\chapter{Sprint S00 --- Initialisation du projet}\label{ch:s00}
\section{Fiche de sprint}
\begin{table}[htbp]\centering\small
\caption{Cadrage du sprint S00}
\begin{tabularx}{\textwidth}{@{}P{31mm}Y@{}}\toprule
\textbf{Champ} & \textbf{Description}\\\midrule
Objectif & Créer un nouveau dépôt Git et un socle technique propre avec une page d’accueil responsive et une API de santé.\\
Périmètre & Spring Boot, Angular, Git, configuration locale, README, architecture simple ; pas de base de données ni d’authentification métier.\\
Stories & \texttt{US-001, US-002, US-003}.\\
Résultat attendu & Backend et frontend qui démarrent ; contenu descriptif sans fausses annonces.\\
Préconditions & Environnement Java/Node/Git vérifié, compatibilité Angular et Node confirmée.\\
Recette & GET santé en HTTP 200 ; frontend compilé ; page responsive ; README reproductible ; dépôt sans secrets.\\
Statut & \tagstatus{PLANIFIÉ} --- démarrage après validation du cadrage.\\
\bottomrule\end{tabularx}
\end{table}

\section{Diagramme de cas d’utilisation du sprint}
\begin{figure}[H]\centering\begin{tikzpicture}[x=1cm,y=1cm,font=\sffamily\small]
\draw[draw=LineGray,rounded corners=2mm] (1,-3) rectangle (11.4,1.1);
\node[anchor=west,font=\sffamily\bfseries\small,text=ForestDark] at (1.4,.7) {Sprint S00 --- socle fonctionnel minimal};
\node[ellipse,draw=Forest,fill=Sage,minimum width=48mm,minimum height=9mm] (ac) at (6.1,-.55) {Consulter la page d’accueil};
\node[ellipse,draw=Forest,fill=Sage,minimum width=48mm,minimum height=9mm] (he) at (6.1,-2.05) {Vérifier la santé de l’API};
\foreach \xx/\yy/\nom in {0/-.6/Visiteur,12.2/-2.0/Développeur}{
 \draw[thick,Forest] (\xx,\yy+.3) circle (.15);
 \draw[thick,Forest] (\xx,\yy+.15) -- (\xx,\yy-.3);
 \draw[thick,Forest] (\xx-.25,\yy-.04) -- (\xx+.25,\yy-.04);
 \draw[thick,Forest] (\xx,\yy-.3) -- (\xx-.22,\yy-.54);
 \draw[thick,Forest] (\xx,\yy-.3) -- (\xx+.22,\yy-.54);
 \node[anchor=north,font=\scriptsize] at (\xx,\yy-.62) {\nom};
}
\draw[Muted] (.3,-.55) -- (ac.west);
\draw[Muted] (11.9,-2.05) -- (he.east);
\end{tikzpicture}

\caption{S00 --- diagramme de cas d’utilisation}\label{fig:s00-uc}\end{figure}
L’initialisation technique n’est pas un parcours commercial ; les deux cas du sprint concernent la \textbf{présentation du service} et la \textbf{vérification de disponibilité}. Les cas de catalogue persistant commencent au sprint S01.

\section{UC-S00-01 --- Consulter la page d’accueil}
\begin{longtable}{@{}P{35mm}P{108mm}@{}}\caption{Description du cas UC-S00-01}\label{tab:s00-uc1}\\
\toprule\textbf{Rubrique} & \textbf{Description}\\\midrule\endfirsthead\toprule\textbf{Rubrique} & \textbf{Description}\\\midrule\endhead
\UCrow{Code}{UC-S00-01}
\UCrow{Acteur principal}{Visiteur}
\UCrow{Objectif}{Lire la présentation de Souk El Hayawan et comprendre les services envisagés.}
\UCrow{Précondition}{Le frontend est lancé et son URL est accessible.}
\UCrow{Déclencheur}{Le visiteur ouvre la page d’accueil.}
\UCrow{Scénario nominal}{1. Le navigateur charge l’application.\newline 2. Angular initialise la page.\newline 3. Le visiteur consulte l’identité, le périmètre et les catégories de services.}
\UCrow{Alternative}{En cas de chargement impossible, le navigateur présente un état d’erreur.}
\UCrow{Postcondition}{La page descriptive est affichée sans résultat de recherche fictif.}
\UCrow{Critère de recette}{Affichage desktop/mobile sans débordement ; pas de bouton prétendant appeler une fonctionnalité inexistante.}
\bottomrule
\end{longtable}

\begin{figure}[H]\centering\begin{tikzpicture}[x=1cm,y=1cm,font=\sffamily\scriptsize]
\foreach \x/\label in {0.5/Visiteur,3.3/Navigateur,6.2/Angular,9.2/{Page accueil}}{
 \node[draw=Forest,rounded corners=1mm,fill=Sage,minimum width=21mm,minimum height=8mm,align=center] at (\x,0) {\label};
 \draw[densely dashed,Muted] (\x,-.4) -- (\x,-4.15);
}
\draw[-{Stealth},Forest] (.5,-1) -- node[above]{Saisir l’URL} (3.3,-1);
\draw[-{Stealth},Forest] (3.3,-1.65) -- node[above]{Charger application} (6.2,-1.65);
\draw[-{Stealth},Forest] (6.2,-2.3) -- node[above]{Construire la vue} (9.2,-2.3);
\draw[-{Stealth},Muted,dashed] (9.2,-2.95) -- node[above]{Contenu statique} (6.2,-2.95);
\draw[-{Stealth},Muted,dashed] (6.2,-3.58) -- node[above]{Afficher accueil} (.5,-3.58);
\end{tikzpicture}

\caption{UC-S00-01 --- diagramme de séquence de consultation}\label{fig:s00-seq-accueil}\end{figure}

\section{UC-S00-02 --- Vérifier la disponibilité de l’API}
\begin{longtable}{@{}P{35mm}P{108mm}@{}}\caption{Description du cas UC-S00-02}\label{tab:s00-uc2}\\
\toprule\textbf{Rubrique} & \textbf{Description}\\\midrule\endfirsthead\toprule\textbf{Rubrique} & \textbf{Description}\\\midrule\endhead
\UCrow{Code}{UC-S00-02}
\UCrow{Acteur principal}{Développeur}
\UCrow{Objectif}{Vérifier que Spring Boot est accessible via une route de santé.}
\UCrow{Précondition}{Le serveur backend est démarré et son port est connu.}
\UCrow{Déclencheur}{Envoi d’une requête GET vers l’endpoint de santé.}
\UCrow{Scénario nominal}{1. Le développeur envoie la requête.\newline 2. Le contrôleur construit un statut.\newline 3. Le client reçoit le code HTTP 200 et un JSON exploitable.}
\UCrow{Alternative}{Backend indisponible : erreur de connexion ou réponse non conforme.}
\UCrow{Postcondition}{L’état de disponibilité est objectivement vérifié.}
\UCrow{Critère de recette}{Code HTTP 200, contrat de réponse documenté et commande de test reproductible.}
\bottomrule\end{longtable}
\begin{figure}[H]\centering\begin{tikzpicture}[x=1cm,y=1cm,font=\sffamily\scriptsize]
\foreach \x/\label in {.5/Développeur,3.6/{Client HTTP},6.8/{HealthController},10.1/{Réponse API}}{
 \node[draw=Forest,rounded corners=1mm,fill=Sage,minimum width=23mm,minimum height=8mm,align=center] at (\x,0) {\label};
 \draw[densely dashed,Muted] (\x,-.4) -- (\x,-4.2);
}
\draw[-{Stealth},Forest] (.5,-1) -- node[above]{Exécuter test} (3.6,-1);
\draw[-{Stealth},Forest] (3.6,-1.72) -- node[above]{GET /api/health} (6.8,-1.72);
\draw[-{Stealth},Forest] (6.8,-2.44) -- node[above]{Construire état UP} (10.1,-2.44);
\draw[-{Stealth},Muted,dashed] (10.1,-3.16) -- node[above]{HTTP 200 / JSON} (3.6,-3.16);
\draw[-{Stealth},Muted,dashed] (3.6,-3.88) -- node[above]{Résultat vérifiable} (.5,-3.88);
\end{tikzpicture}

\caption{UC-S00-02 --- diagramme de séquence de vérification API}\label{fig:s00-seq-health}\end{figure}

\section{Diagramme de classes de conception}
\begin{figure}[H]\centering\begin{tikzpicture}[font=\sffamily\small]
\node[draw=Forest,rectangle split,rectangle split parts=3,rounded corners=1mm,fill=Sage,text width=46mm,align=left] (c) at (0,0) {\textbf{HealthController}\nodepart{second}\texttt{- service: HealthService}\nodepart{third}\texttt{+ health(): HealthStatus}};
\node[draw=Forest,rectangle split,rectangle split parts=3,rounded corners=1mm,fill=Sage,text width=48mm,align=left] (s) at (7,0) {\textbf{HealthStatus}\nodepart{second}\texttt{+ status: String}\nodepart{third}\textit{Objet de réponse HTTP}};
\draw[-{Stealth[length=2mm]},thick,Forest] (c.east) -- node[above,font=\scriptsize]{retourne} (s.west);
\node[font=\scriptsize,text=Muted,align=center] at (3.5,-2) {Modèle pédagogique de conception : aucun objet métier ni couche JPA au Sprint 0};
\end{tikzpicture}

\caption{S00 --- premier diagramme de classes technique, sans données métier}\label{fig:s00-class}\end{figure}

\section{Diagramme d’activité}
\begin{figure}[H]\centering\begin{tikzpicture}[node distance=5mm and 8mm,font=\sffamily\scriptsize,>=Stealth]
\tikzset{act/.style={draw=Forest,fill=Sage,rounded corners=2mm,minimum width=42mm,minimum height=8mm,align=center}}
\node[circle,fill=ForestDark,minimum size=5mm,inner sep=0pt] (start) {};
\node[act,below=of start] (open) {Ouvrir l’URL du site};
\node[diamond,draw=Forest,fill=Warm,aspect=2,minimum width=31mm,minimum height=11mm,align=center,below=6mm of open] (ok) {Application accessible ?};
\node[act,below=8mm of ok] (view) {Afficher l’accueil};
\node[act,right=15mm of ok] (err) {Afficher un état d’erreur};
\node[draw=ForestDark,circle,double,double distance=1pt,minimum size=5mm,inner sep=0pt,below=of view] (end) {};
\draw[->,Forest] (start)--(open);
\draw[->,Forest] (open)--(ok);
\draw[->,Forest] (ok)--node[right]{Oui}(view);
\draw[->,Forest] (ok)--node[above]{Non}(err);
\draw[->,Forest] (view)--(end);
\draw[->,Muted] (err.south) |- (end.east);
\end{tikzpicture}

\caption{S00 --- activité de consultation de l’accueil}\label{fig:s00-act}\end{figure}

\section{Maquette de l’accueil --- style Balsamiq}
\begin{figure}[H]\centering\resizebox{\textwidth}{!}{\begin{tikzpicture}[x=1cm,y=1cm,font=\sffamily\scriptsize]
\fill[LightGray] (0,0) rectangle (13.5,-7.2);
\draw[Muted,thick,rounded corners=1mm] (0,0) rectangle (13.5,-7.2);
\fill[gray!22] (0,0) rectangle (13.5,-.72);
\node[anchor=west,font=\bfseries] at (.35,-.35) {SOUK EL HAYAWAN};
\foreach \x/\t in {8.9/Catalogue,10.5/Entraide,12.0/Connexion}{\node at (\x,-.35) {\t};}
\draw[densely dashed,gray!70] (.45,-1.15) rectangle (13,-3.45);
\node[font=\normalsize\bfseries] at (6.75,-1.8) {Tous les animaux, une même plateforme};
\node[font=\scriptsize,text=Muted,align=center] at (6.75,-2.3) {Vente, adoption, don et entraide --- présentation du projet};
\fill[gray!40] (5.85,-2.68) rectangle (7.65,-3.13);
\node at (6.75,-2.9) {Découvrir};
\node[anchor=west,font=\small\bfseries] at (.55,-4.02) {Les services proposés};
\foreach \x/\t/\tB in {.65/Vente/{Annonce},4.82/Adoption/{Gratuite},8.99/Entraide/{Disparus}}{
 \draw[gray!70,rounded corners=1mm,fill=white] (\x,-4.35) rectangle ++(3.75,-1.8);
 \fill[gray!24] (\x+.24,-4.6) rectangle ++(3.25,-.65);
 \draw[gray!50] (\x+.24,-4.6) -- ++(3.25,-.65);
 \draw[gray!50] (\x+3.49,-4.6) -- ++(-3.25,-.65);
 \node[font=\bfseries] at (\x+1.875,-5.55) {\t};
 \node[text=Muted] at (\x+1.875,-5.86) {\tB};
}
\node[anchor=west,text=Muted] at (.5,-6.75) {À propos \quad | \quad Contact \quad | \quad Conditions d’utilisation};
\end{tikzpicture}
}
\caption{S00 --- maquette basse fidélité de la landing page (illustration)}\label{fig:s00-wire}\end{figure}
\begin{notebox}
\textbf{Décision de conception :} S00 ne présente pas de faux comptes ou de fausses annonces. La maquette représente une page \emph{informative}. Les écrans de catalogue reliés à la base réelle interviennent en S01.
\end{notebox}

\section{Travaux du développeur et validation}
\begin{enumerate}
\item Rédiger trois parcours utilisateur avec ses propres mots.
\item Créer le dépôt Git, organiser \texttt{frontend/} et \texttt{backend/}, configurer \texttt{.gitignore}.
\item Initialiser Spring Boot et créer l’endpoint de santé.
\item Initialiser Angular, SCSS et les composants de base.
\item Réaliser l’accueil, vérifier le responsive et les liens actifs.
\item Produire un README et une démonstration de fonctionnement.
\end{enumerate}
\appendixplaceholder{résultats des tests, captures de l’application développée, détails d’implémentation et bilan personnel du sprint S00}
